\section{犯罪构成 - 定罪标准：两阶层的犯罪构成体系}
\begin{enumerate}
  \item 事实判断与价值评价。实然/应然。
        \begin{enumerate}
          \item 事实判断。\zd{客观标准、真假检测}。\al{下暴雨 vs 天气真糟糕、腿短 vs 裙子不好看、男 vs 渣男}。
        \end{enumerate}
  \item 定罪步骤。事实判断 \yjt 价值评价 \yjt 构成犯罪 \yjt 负刑事责任（法律后果）。
        \begin{enumerate}
          \item 事实判断。\zd{法益侵害事实}/违法事实 \yjt \zd{法益侵害性}/违法性 \yjt 客观违法阶层（违法性）。\al{不是我干的}。
          \item 价值评价。刑法能否谴责你 \yjt 主观责任阶层（可谴责性 \bdy 负刑事责任）。\al{我不是故意的}。
        \end{enumerate}
  \item \bs{两阶层体系}。
        \begin{enumerate}
          \item 客观要件。
                \begin{enumerate}
                  \item \zx{行为主体、危害行为、危害结果、因果关系} \yjt 初步具有违法性。立。
                  \item 客观阻却事由。\zx{正当防卫、紧急避险、（推定的）被害人承诺}、\scx{义务冲突}。破。
                \end{enumerate}
          \item 主观要件。
                \begin{enumerate}
                  \item \zx{故意、过失}、\zd{无罪过事件}（\yi 不可抗力。\er 意外事件）、\zd{事实认识错误} \yjt 能够谴责。立。
                  \item 主观阻却事由。\zx{责任年龄、责任能力、违法性认识可能性}、\bs{期待可能性}。破。
                \end{enumerate}
        \end{enumerate}
  \item 应用顺序。客观主义立场。\bs{先客观，后主观}！思想不能犯罪。\zx{不能随意排列组合}。
  \item \zd{鉴定式案例分析法}。按照体系顺序去审查每一个罪的构成要件。
  \item 三阶层。构成要件\zx{该当性、违法性、有责性}。太复杂，不适合实务。
  \item 四要件的 bug。要求犯罪主体具有责任年龄、责任能力。放到共犯一起考。\al{10 岁甲要强奸乙，让爷爷望风}。
\end{enumerate}
